    \section{Resume}

    % skal være i kursiv

    \textit{
    Dette projekt har til hensigt at besvare, hvordan man kan hylde fejringen af Viborg Katedralskoles 100-års jubilæum for bygningen ved hjælp af et spil. 
    
    Med udgangspunkt i den historiske begivenhed den 23. oktober 1944, hvor Gestapo gennemførte en razzia på skolen og arresterede 14 modstandsfolk, udvikles en digital oplevelse, der formidler skolens besættelseshistorie til nuværende og kommende elever i aldersgruppen 14-19 år. 
    
    Med henblik på en større detaljegrad samt ønsket om en nøjagtig repræsentation af skolen blev photogrammetry anvendt til skabelsen af realistiske 3D-modeller af skolens omgivelser og lokaler. 

    Projektet anvender Scrum og Unified Process som styringsmetoder og udvikles i Unreal Engine med fokus på multiplayer-elementer. 
    
    Spillets systemer er defineret gennem analyse- og designfaser som identificerer kerne aspekter af spillet og uddyber deres funktionalitet. 

    Spillet blev udviklet i Unreal Engine ved brug af Blueprints og C++ som sammenlagt håndterer henholdsvis gameplay-logik og UI interaction samt mere komplekse implementationer.  Spillet fungerer ved brug af en klient-server-arkitektur og Unreal Engines replication-system. 

    Resultatet er en spil-prototype som har multiplayer funktionalitet, men som ikke nødvendigvis reflekterer det historiske perspektiv af bygningen, som projektet bygger på. 
    }